\section{One Borel graph rule}
\label{sec:implementation}

We implement the preceding construction using only vertex labels and the
distinguished undirected edge. Local finiteness and connectivity imply that
each graph has countably many vertices and edges.

\subsection{From vertex labels to independent edge paths}

\begin{lemma}\label{lem:edge-noise}
There is a Borel equivariant rule assigning a continuous path $w_e$ starting
at zero to each edge of a vertex-labeled graph such that, for every fixed
graph with independent uniform vertex labels, $(w_e)_{e\in E(G)}$ are
independent standard Brownian motions.
\end{lemma}

\begin{proof}
Partition the binary digits of a uniform variable into countably many
infinite subsequences. With a fixed Borel convention at dyadic rationals,
this gives a Borel map from one uniform variable to a countable family of
independent uniforms. At each vertex $v$, use this map to obtain a key $K_v$
and a private sequence $(V_{v,r})_{r\ge1}$, with all these variables independent.
Almost surely the keys are distinct throughout the graph. On this event,
an edge $e=\{u,v\}$ with $K_u<K_v$ selects $V_{u,r}$, where $r$ is the rank
of $v$ among the neighbors of $u$, ordered by key. The rank is finite.
Different edges select different pairs $(u,r)$: edges with different source
vertices use different pools, and edges with the same source vertex have
different other endpoints and hence different ranks. Simplicity is used here.
Conditional on all the keys, the selected variables are distinct coordinates
of an independent uniform family that is independent of those keys. They
therefore have the product uniform law.

For completeness, one uniform can be mapped Borel measurably to a standard
Brownian path. Split it into uniforms and apply the normal quantile to
obtain independent standard normals, using finite defaults at exceptional
quantile inputs. Start at $w(0)=0$ and use some of them as independent increments on the
nonnegative integer grid. On an interval $[a,a+h]$ of a dyadic grid, insert
the midpoint value
\[
 \frac{w(a)+w(a+h)}2+\frac{\sqrt h}{2}\,\xi,
\]
where $\xi$ is a fresh independent standard normal. The two new increments
are jointly Gaussian, have variance $h/2$, and have covariance zero. By
induction, every dyadic grid has independent Gaussian increments of the
appropriate variances.

On $[0,K]$, with $K$ a positive integer, the refinement of mesh $2^{-p}$
uses $K2^p$ new normals. The probability that any has absolute value greater
than $p$ is at most $2K2^p e^{-p^2/2}$, which is summable for $p\ge1$.
The change of linear interpolation at that level is therefore eventually
at most $p2^{-p/2}/2$. These bounds are summable. The interpolations converge
uniformly on every compact interval almost surely. Their limit has continuous
paths and the required dyadic Gaussian laws, hence the Wiener law by
continuity. The event of compact-uniform convergence is Borel: it is the
countable collection of uniform Cauchy conditions on integer intervals.
Each condition can be tested at rational times for these continuous
interpolations. Taking the limit on this event and the zero path otherwise
defines a Borel path-valued map on every input.

Apply this same map to every selected edge uniform. This produces the
required product Wiener law. On a labeling with any repeated keys, assign
the zero path to every edge. A key collision is witnessed in some finite
ball about the distinguished edge, so this exceptional case is Borel and
equivariant. It has probability zero on each fixed countable graph. All
other selections involve only finite neighbor lists and are Borel and
equivariant as well.
\end{proof}

\subsection{Measurability on varying graphs}

The decoder must be one rule on graph space, in addition to being Borel for
each fixed graph. We give the details of this point. On the Borel set of
labelings with distinct keys, vertices in any finite ball can be listed in
increasing key order. Its undirected edges can then be listed lexicographically by
$(\min\{K_u,K_v\},\max\{K_u,K_v\})$. These are finite measurable listings, not a global ordering
of the vertices. Re-distinguishing the graph at any edge selected from such
a list is a Borel operation: each finite neighborhood of that edge can be
read in a sufficiently large finite neighborhood of the original edge.
For example, if the selected edge lies in $H_r(e)$, its radius-$q$
neighborhood is visible in $H_{r+q+1}(e)$. In particular, a Borel edge rule
may be evaluated measurably at the finitely many edges in any such ball.
The evaluated rule may itself depend on the entire graph; measurability
follows by composition with re-distinguishing, not from finite dependence.

For the drift from Section~\ref{sec:response}, every
$b_{e,S_n(e)}(t,y)$ is the limit of a finite weighted-tree ratio on $H_m(e)$
as $m\to\infty$, with $n$ held fixed. The finite ratio uses only finitely
many current fields and the finite graph. The outer limsup in $n$ then
defines $b_e(t,y)$. This exact order of limits makes the formula Borel when
the current fields are supplied by any Borel edge rule jointly measurable
in $t$.

Start with the paths $Z_e^0(t)=w_e(t)$ from Lemma~\ref{lem:edge-noise}.
They are jointly Borel in the labeled graph, the distinguished edge, and
$t$. If the same is true of $Z^k$, the preceding finite listings and drift
formula show that
\[
 (G,U,e,s)\longmapsto b_e(s,Z^k(s))
\]
is Borel and bounded. Its scalar parameter integral over $[0,t]$ is Borel
in $(G,U,e,t)$: on $0\le t\le T$ it is the integral over $[0,T]$ of the
bounded Borel function multiplied by $\1_{\{s\le t\}}$. This is the usual
measurability of integration with respect to a fixed finite measure, proved
first for indicator rectangles and then by bounded approximation. Thus
the next Picard iterate is jointly Borel. The countable limit defining
$Z$, and the integer-time limsup in \eqref{eq:readout}, preserve Borel
measurability.

Every operation commutes with labeled isomorphisms. In particular, the balls
$H_n(e)$ are centered on the unordered endpoints of $e$, and the finite
key order is transported by any labeled isomorphism. Scalar time integration
and the prescribed countable limits preserve this invariance. There is no
distinguished vertex or external enumeration in the construction.

\begin{proof}[Proof of Theorem~\ref{thm:main}]
On labelings with distinct keys, define $\Phi(G,U,e)$ by the readout
\eqref{eq:readout}, using the intrinsic forest drift and the edge paths of
Lemma~\ref{lem:edge-noise}. On the remaining labelings set $\Phi=0$ at every
edge. The preceding argument makes this a single Borel, root-independent,
equivariant rule, defined on all labelings.

Fix a graph $G$. Its key-collision event has probability zero, and the edge
paths have product Wiener law. Lemma~\ref{lem:forest-response} gives the
uniform finite-tilt response bound with $L=1/2$ for $\nu=\FUSF_G$.
Proposition~\ref{prop:brownian-sampler} therefore identifies the joint law
of all the output indicators as $\FUSF_G$. This proves the assertion for
every fixed graph.
\end{proof}

\begin{proof}[Proof of Corollary~\ref{cor:unimodular}]
Place independent uniform labels conditionally on $(G,o)$. The graphwise
conclusion of Theorem~\ref{thm:main}, for the same Borel rule $\Phi$, gives
the desired conditional law. This is an equality of conditional probability
kernels for each graph; it requires no intersection of label-conull sets
over different graphs. Unimodularity imposes no further requirement on the
construction.
\end{proof}
